

  \documentclass[12pt, german, english]{article}
    \usepackage[T1]{fontenc}
    \usepackage[ansinew]{inputenc}
    \usepackage{babel}
    \usepackage{fullpage}
 %   \usepackage{english}

\begin{document}
\title{\bf  Haavelmo's theorem in open economies: the insulation effect}
\author{Nikolaus K.A. L�ufer\\
Professor em., University of Konstanz, Germany} \vspace{-1.5mm}
\date{25.5.2009}
\maketitle

\abstract{The Haavelmo-Theorem is always treated in a closed economy
context. What happens to the theorem in an open economy? In the
present paper this issue is treated for two cases: one where imports
depend on disposable income, and in the other case imports are
assumed to depend on total income. For the first case (disposable
income dependency of imports), the theorem may be confirmed, for the
second case (total income dependency of imports) the multiplier
becomes even smaller than one. In the first case, the domestic
fiscal expansion does not spill over to the foreign country, the two
countries are insulated. In a forthcoming paper I shall demonstrate
that disposable income dependency of imports is the preferable
alternative (case) on a priori theoretical grounds.}

\section{The Haavelmo-Theorem in an open economy}
\subsection{Introduction}
The Haavelmo theorem specifies the income changes following a tax
financed expansion of public expenditures. This income change is
expressed by a multiplier applied to the change in public
expenditures. Specifically, a multiplier of exactly 1 is determined
by the theorem. In textbooks, the Haavelmo-Theorem arises in a
closed economy context. The consequences of opening up the economy
are never investigated.

The present note attempts to show, that the Haavelmo result of a
unity multiplier remains intact when the economy is opened up.
Whether exports in the home country are exogenous or endogenous does
not matter for the result. What really matters is the kind of income
dependency of imports in the home country. These imports must depend
on private disposable income in the same way as does private
consumption. If this assumption is replaced by an import function
that depends on income instead of disposable income, then the
multiplier turns out to be smaller than 1.


\subsection{Endogenous exports and imports in a
two-country setting} We first consider the case of endogenous
exports of the home country and start with the income equation of
the home country in a two-country world.

\begin{equation}
 Y_1 = C(Y_1-T) + I + A + Ex(Y_2) - Im(Y_1-T)
\end{equation}

\noindent The symbols have the usual meaning as known from
macroeconomics. A stands for state expenditures, T for taxes, Y for
income, I for investments, Ex for exports and Im for imports. For
our demonstration, we need the differential form of this
equation.
\begin{equation} dY_1 = c_1 (dY_1-dT) + dI + dA + m_2 dY_2
- m_1(dY_1-dT).
\end{equation}
Here, $c_i$ and  $m_i$  represent, respectively, the marginal rates
of consumption and of imports of country i (i=1,2).

We consider the case of $dI=0$ and assume, as is standard in the
Haavelmo-Theorem for closed economies:
\begin{equation}
dA = dT.
\end{equation}

$dY_2$ is determined by means of the differential of the income
equation of the foreign country:
\begin{equation}
dY_2 = c_2 dY_2 + m_1(dY_1-dT) - m_2dY_2.
\end{equation}
The term $m_1(dY_1-dT)$ represents the additional exports of the
foreign country to the home country. Solving this equation for
$dY_2$, we obtain:
\begin{equation}
dY_2 =\frac{m_1}{1-c_2+m_2}(dY_1-dT).
\end{equation}

Inserting this expression for $dY_2$ into equation (2) and solving
for $dY_1$, and observing (3), one arrives at:
\begin{equation}
dY_1 = dA,
\end{equation}
which implies a multiplier of 1.

\subsection{Exogenous exports and endogenous imports} The
case of exogenous exports may be treated in the chosen framework by
setting $m_2 dY_2=0$ in equation (2). The result of a unitary
multiplier remains intact.



\section{The insulation effect}
\subsection{Spill-over effects to the foreign country}
The question we shall address next is, whether there are spill-overs
from the home country  to the foreign country linked with tax
financed public expenditure extensions of the home country.

\noindent Inserting the result, as expressed by equation (6), into
equation (5) and using (3), we obtain
\begin{equation}
dY_2 = 0.
\end{equation}

\noindent This equation tells us that a tax financed change in
public expenditures in the home country, $dY_1=dA = dT \neq 0$, does
not spill-over to the foreign country. The foreign country is
insulated from tax financed expansions of government expenditures in
the home country.

%In an international context of two countries, the
%Haavelmo-theorem extends to an insulation theorem and the following
%statements hold: Tax financed extensions of public expenditures do
%not spill-over to the foreign country. The foreign country is
%insulated from tax financed extensions of public expenditures in the
%home country.

\subsection{Repercussions from the foreign country}
Repercussion effects, i.e. effects that run from the foreign country
back to the home country are considered next.

Since tax financed extensions of public expenditures do not
spill-over to the foreign country, international repercussions
cannot occur. The home country is protected against international
repercussions by tax financing its extensions of public
expenditures.

%\subsection{Can insulation be avoided in the long run?}
%Quite the opposite. Since tax financing of public expenditures
%cannot be avoided in the long run, insulation must be obtained in
%the long run.

\subsection{Economic interpretation} With complete tax financing of
additional public expenditures, an increase of public expenditures
does not increase private disposable income. Consequently, imports
will not rise. The trickling away of demand by imports is zero and
from the foreign country no repercussions will result even if
exports of the home country are not treated, from the start, as
outright exogenous. Can insulation be avoided in the long run? Quite
the opposite. Since tax financing of public expenditures cannot be
avoided in the long run, insulation must be obtained in the long
run.

\subsection{An insulation theorem} The unitary multiplier of the
Haavelmo theorem for a closed economy holds also for the case of an
open economy. However, in an open economy, tax financing of public
expenditures also implies an international insulation effect.
Spill-over effects from the home country to foreign countries and
international repercussions from foreign countries to the home
country are absent. This result may form the content of an
insulation theorem.

\section{Alternating the import function}
Until now, the import function of the home country depended on
disposable income. If, instead, we suppose dependence on (total)
income then we have to replace the term $m_1(dY_1-dT)$ in equations
(2) and (4) by $m_1dY_1$. With
\begin{equation}
dY_2 = \frac{m_1}{1-c_1+m_2} dY_1,
\end{equation}
for the change of total income we obtain
\begin{equation}
dY_1 =\frac{(1-c_1)(1-c_2) +(1-c_1)m_2}{ (1-c_1)(1-c_2) +(1-c_2)m_1+
(1-c_1)m_2} dA
\end{equation}
in the case of endogenous exports and
\begin{equation}
dY_1 =\frac{(1-c_1)}{ (1-c_1) + m_1} dA
\end{equation}
in the case of exogenous exports of the home country. In both cases
the multiplier is clearly below 1 and the former insulation result
does not hold anymore. The foreign country receives spill-overs and
there are positive repercussions from the foreign country to the
home country. The positive repercussions can only mitigate but not
reverse the negative effects on home country (total) income
following the leaking of demand by income dependent imports. With a
multiplier smaller than 1, disposable income falls, i.e. $dY_1-dT<0$.




\section{Completing the model and final result}
At this point, we seem to have a choice between a multiplier of 1
and an insulation effect on one side and a multiplier of less than 1
and no insulation on the other side, this choice being linked to the
kind of import dependency assumed, disposable income in the first
case and total income in the second case. In a forthcoming paper, I
shall demonstrate that the import function should depend on
disposable income and not on total income. However, the models
considered so far do not include the financial sector apart from the
requirement to finance expenditures by taxes. In particular, the
underlying behavioural functions do not depend on the rate of
interest. Equally, exchange rates have been excluded from
consideration which means that either fixed exchange rates are
assumed or it is assumed that flexible exchange rates, similar to
interest rates, do not affect behaviour.

Now, if dependency on interest rates and on exchange rates are
explicitly introduced, as long as the fixed price assumption for
commodities and wages is still maintained, any constraints coming
from the side of the financial sector in the form of a less than
perfect elasticity of the supply of money will  reduce the
multiplier even further below 1. A flexible exchange rate may
attenuate any lowering of the multiplier but will not be able to
reverse the result of a multiplier smaller than 1 into a multiplier
not smaller than 1.

Thus, the Haavelmo theorem with its unitary multiplier is
establishing an upper boundary for the size of the multipliers under
less limiting assumptions. It is definitely not a lower boundary, as
some have maintained.



\bibliographystyle{plain}
\begin{thebibliography}{}
T. Haavelmo, Multiplier effects of a Balanced Budget, Econometrica,
Vol. 13, 1945, S. 311ff.
\end{thebibliography}

\end{document}
